\documentclass[11pt]{article}
\usepackage[margin=1in]{geometry}
\usepackage[T1]{fontenc}
\usepackage{lmodern}
\usepackage{amsmath,booktabs,hyperref,microtype}
\hypersetup{hidelinks}
\makeatletter\let\tableinput\@@input\makeatother
% Generated; do not edit. Within-model, interval-robust.
\newcommand{\CertAttackVersion}{17.1}
\newcommand{\CertBackupWorstPct}{22.6}
\newcommand{\CertCoreWorstPct}{3.4}
\newcommand{\CertEmptyWorstPct}{28.2}
\newcommand{\CertIdentityWorstPct}{3.4}
\newcommand{\CertMinGuardFive}{3}
\newcommand{\CertMinGuardNarrow}{2}
\newcommand{\CertMinGuardOne}{6}
\newcommand{\CertMinGuardTen}{1}
\newcommand{\CertMinGuardWide}{10}
\newcommand{\CertMitigations}{44}
\newcommand{\CertNamedGuardFive}{2}
\newcommand{\CertNamedTotal}{6}
\newcommand{\CertTechniques}{588}
\newcommand{\CertWideEmptyPct}{28.2}


\title{Distribution-free control-adequacy certificates on the MITRE ATT\&CK kill chain}
\author{}
\date{2026-09-20 (working note, version 0.1.0)}

\begin{document}
\maketitle

\begin{abstract}
We ask a decision question with a real, curated attack model rather than a
synthetic one: which portfolios of defensive controls \emph{provably} keep a
ransomware operation's reachability to its impact objective below a chosen level,
when the effectiveness of each control is only known within an interval? Using the
pinned MITRE ATT\&CK Enterprise knowledge base (v\CertAttackVersion; \CertTechniques{}
techniques, \CertMitigations{} mitigations, and their real \texttt{mitigates}
edges), we define a monotone kill-chain reachability function and certify
portfolios by evaluating it at the adverse and favourable corners of the
effectiveness uncertainty set --- the necessary/possible interval-efficiency
argument, applied to attacker reachability. The guarantee is distribution-free and
holds for every parameter value in the set. We report a cost-versus-guarantee
frontier, a set of realistic named bundles, and how the minimum certifiable
portfolio size grows with assumed uncertainty. The contribution is the method and
its application to real ATT\&CK structure; effectiveness intervals are analyst
priors, not incident measurements, and no claim is made about any specific
organisation.
\end{abstract}

\section{Model}
A ransomware operation traverses the ordered post-compromise ATT\&CK tactics from
initial access to exfiltration and then the impact objective T1486 (Data Encrypted
for Impact). A control portfolio is a set of ATT\&CK mitigations. Each technique's
residual per-attempt success is $\mathrm{base}\cdot\prod_{m}(1-\mathrm{eff}_m)$ over
the portfolio mitigations that mitigate it. A stage's success is the
usage-weighted mean residual over its techniques (weights are the real ATT\&CK
procedure-usage counts, so a prevalent technique matters more and partial coverage
helps proportionally); the impact stage contributes the objective technique alone.
Reachability $R$ is the product of stage successes. $R$ is coordinatewise monotone:
increasing in every base rate and decreasing in every $\mathrm{eff}_m$.

\section{Certificate}
Effectiveness is uncertain, so each $\mathrm{eff}_m\in[e^{\mathrm{lo}}_m,e^{\mathrm{hi}}_m]$
and the base rate lies in a common interval. Monotonicity places the extremes of
$R$ at the box corners. A portfolio is \emph{guaranteed} adequate at level
$\varepsilon$ when $R\le\varepsilon$ at the worst corner (base high, every
effectiveness low), and only \emph{possible} when $R\le\varepsilon$ merely at the
best corner. This is exactly the necessary/possible efficiency corner argument used
elsewhere in this repository, here applied to reachability; where incident counts
exist, finite-sample betting bounds supply the intervals.

\section{Results}
With a common base-rate prior $[0.5,0.9]$ and a wide effectiveness prior
$[0.20,0.70]$ (multi-factor authentication tightened to $[0.85,0.99]$ from the
reported account-compromise reduction), the undefended worst-case reachability is
\CertEmptyWorstPct\%. Along the greedy cost-versus-guarantee frontier,
\CertMinGuardTen{} mitigation guarantees $R\le 10\%$, \CertMinGuardFive{} guarantee
$R\le 5\%$, and \CertMinGuardOne{} guarantee $R\le 1\%$ (Table~\ref{tab:frontier}).
Among \CertNamedTotal{} realistic named bundles, \CertNamedGuardFive{} are
guaranteed adequate at $5\%$: the identity bundle (worst case
\CertIdentityWorstPct\%) and the hospital-core bundle (\CertCoreWorstPct\%), whereas
data backup alone leaves \CertBackupWorstPct\% (Table~\ref{tab:named}). Identity
controls are the dominant lever; single controls from other families do not certify
even $R\le 10\%$.

\begin{table}[t]\centering
\caption{Greedy worst-case-reduction frontier (first eight additions): mitigation
added and worst-case reachability after it, percent.}\label{tab:frontier}
\begin{tabular}{@{}rllr@{}}\toprule
Size & Mitigation & Name & Worst $R$ (\%)\\\midrule
\tableinput table_control_frontier_rows.tex
\bottomrule\end{tabular}
\end{table}

\begin{table}[t]\centering
\caption{Named control bundles: size, worst-case reachability, and adequacy
certificates.}\label{tab:named}
\begin{tabular}{@{}lrrll@{}}\toprule
Bundle & Size & Worst $R$ (\%) & Guar.\ $\le$10\% & Guar.\ $\le$5\%\\\midrule
\tableinput table_control_named_rows.tex
\bottomrule\end{tabular}
\end{table}

\section{Sensitivity to assumed uncertainty}
The certificate's demands scale with ignorance. To guarantee $R\le 5\%$, a narrow
effectiveness prior $[0.35,0.55]$ needs \CertMinGuardNarrow{} mitigations, while a
wide prior $[0.10,0.85]$ needs \CertMinGuardWide{} (with undefended worst case
\CertWideEmptyPct\%). This is the intended behaviour: a distribution-free guarantee
buys robustness at the price of demanding more coverage when less is known.

\section{Limitations}
ATT\&CK's \texttt{mitigates} edges are qualitative expert mappings, not measured
effect sizes; the effectiveness intervals are analyst priors, with one
evidence-tightened exception, and the kill-chain progression is a modelling
abstraction. There is no validation against real incidents. The result is a
method and its application to real ATT\&CK structure with explicit uncertainty
handling --- a decision-support certificate, not an empirical estimate of any
organisation's risk.

\paragraph{Data and reproducibility.} MITRE ATT\&CK Enterprise v\CertAttackVersion{}
is used unmodified under the ATT\&CK Terms of Use; the pinned bundle is hashed and
all outputs carry a provenance manifest. Every number above is generated from the
committed result tables.

\end{document}
